\chapter{Introducción}\label{cap:introduccion}

A lo largo de los últimos años, la adopción de servicios en la nube se ha consolidado
como el paradigma predominante para el despliegue de aplicaciones y la gestión de
infraestructura. Sin embargo, detrás de la aparente simplicidad que los proveedores
cloud prometen, subyace un conjunto de complejidades técnicas, económicas y
organizacionales que condicionan profundamente las decisiones de quienes trabajan
con estas tecnologías. Este capítulo presenta la motivación que da origen a este
trabajo, los objetivos que se persiguen, la estructura que sigue la memoria y
una declaración explícita sobre el uso de herramientas de inteligencia artificial
durante su elaboración.

\section{Motivación}

Mi interés por las tecnologías cloud surgió de forma gradual durante los años de
formación, a través de la exploración de distintos campos de la ingeniería del
software. Mi acercamiento a \textbf{Amazon Web Services (AWS)} —y la obtención de
la certificación \textbf{AWS Cloud Practitioner}— me permitió comprender de primera
mano tanto las posibilidades que ofrece la nube como las fricciones que aparecen en
cuanto se profundiza en su uso real.

Una de las primeras fricciones que identifiqué fue el \textbf{compromiso tecnológico}
que exige cada proveedor. Aunque los conceptos fundamentales de la computación en
la nube son transferibles entre plataformas, las tecnologías concretas no lo son:
los servicios, las nomenclaturas, las herramientas de gestión y los modelos de
configuración son específicos de cada proveedor. Aprender a desplegar en AWS no
me preparaba para desplegar en \textbf{Microsoft Azure} o en \textbf{Google Cloud
Platform (GCP)}; cada uno exige su propio proceso de aprendizaje.

A esta dificultad se suma la opacidad en la estimación de costes. Aunque los
proveedores ofrecen herramientas de cálculo —como la calculadora de costes de
AWS—, su uso requiere conocimiento previo de los servicios contratados, acceso a
una cuenta activa y criterio suficiente para modelar correctamente el escenario
de despliegue. En mis primeros pasos con la nube me encontré con que estimar qué
iba a costar un despliegue era, en sí mismo, un problema no trivial.

Con el tiempo, la experiencia laboral me permitió observar estas mismas fricciones
a escala organizacional. En entornos profesionales, la elección de un proveedor cloud
no es una decisión puramente técnica: implica estudios económicos comparativos,
negociaciones con proveedores, formación especializada del equipo y, en muchos
casos, la dependencia de empresas externas especializadas en la gestión de
infraestructura. Lo que conceptualmente debería ser sencillo —desplegar una
aplicación en la nube— se convierte en un proceso complejo que requiere recursos
y conocimientos que no siempre están disponibles.

Es en este contexto donde surge la motivación para desarrollar \textbf{CloudPorter},
una herramienta que permita describir infraestructura de forma agnóstica al
proveedor y desplegarla en distintos clouds, reduciendo el coste de entrada,
facilitando la comparación económica entre alternativas y disminuyendo el grado
de compromiso tecnológico que hoy exige la nube.

\section{Objetivos}

Este trabajo no se concibe como la implementación de un sistema cerrado y completo,
sino como una exploración fundamentada de un problema real —el \textbf{cloud lock-in}— con el objetivo de construir una base sólida sobre la
que seguir avanzando. Los requisitos que se presentan en el capítulo~\ref{cap:especificacion}
describen el sistema en su forma ideal; los objetivos que se enuncian aquí definen
la dirección del trabajo y el alcance que se persigue alcanzar.

El \textbf{objetivo general} de este TFG es desarrollar \textbf{CloudPorter}, una
herramienta que permita describir infraestructura cloud de forma agnóstica al
proveedor y desplegarla en distintas plataformas, reduciendo el compromiso
tecnológico y económico que hoy exige trabajar con la nube.

Para guiar el desarrollo hacia ese objetivo, se identifican los siguientes
\textbf{objetivos específicos}:

\begin{enumerate}
    \item \textbf{Definir una abstracción cloud-agnóstica} que permita describir
    infraestructura sin acoplarse a un proveedor concreto.

    \item \textbf{Demostrar que esa abstracción es desplegable} en múltiples
    plataformas cloud reales.

    \item \textbf{Incorporar visibilidad económica} que permita comparar el coste
    entre proveedores.

    \item \textbf{Ofrecer una interfaz accesible} que no exija conocimiento
    profundo de cada proveedor para operar la herramienta.
\end{enumerate}

\section{Estructura de la memoria}

La memoria se organiza en ocho capítulos que siguen la progresión natural del
trabajo: desde la definición del problema hasta el análisis de los resultados
obtenidos y las líneas de trabajo futuro.

\begin{itemize}
    \item \textbf{Capítulo~\ref{cap:descripcion} — Descripción del problema:}
    analiza el cloud lock-in en sus dimensiones técnicas, económicas y
    organizacionales, y presenta la solución propuesta.

    \item \textbf{Capítulo~\ref{cap:estado} — Estado del arte:}
    revisa las herramientas existentes y sitúa CloudPorter en relación con ellas.

    \item \textbf{Capítulo~\ref{cap:especificacion} — Especificación de requisitos:}
    recoge los requisitos del sistema, distinguiendo el alcance del trabajo de
    las líneas de evolución futura.

    \item \textbf{Capítulo~\ref{cap:planificacion} — Planificación y metodología:}
    describe la metodología de desarrollo adoptada, el cronograma y la planificación
    del trabajo.

    \item \textbf{Capítulo~\ref{cap:desarrollo} — Desarrollo:}
    núcleo de la memoria; documenta el desarrollo iterativo de CloudPorter mediante
    ciclos Kanban, con las decisiones de diseño, la implementación y los resultados
    de cada iteración.

    \item \textbf{Capítulo~\ref{cap:evaluacion} — Evaluación:}
    presenta la validación del sistema construido frente a los requisitos definidos.

    \item \textbf{Capítulo~\ref{cap:conclusiones} — Conclusiones:}
    recoge los aprendizajes obtenidos, el alcance alcanzado y las líneas de
    desarrollo futuro que quedan abiertas.
\end{itemize}

\section{Uso de inteligencia artificial}

Durante la elaboración de este Trabajo Fin de Grado se ha recurrido a herramientas de
\textbf{inteligencia artificial (IA)} generativa como apoyo en dos ámbitos diferenciados:
el desarrollo del software CloudPorter y la redacción de la presente memoria. A continuación
se detallan las herramientas utilizadas, los usos concretos realizados y los mecanismos de
validación empleados para garantizar la corrección y coherencia del trabajo entregado.

\subsection{Herramientas utilizadas}

La única herramienta de IA empleada a lo largo del proyecto ha sido \textbf{Claude Code}
\cite{claudeCode}, un agente de programación e interacción en lenguaje natural desarrollado
por Anthropic, que opera desde la línea de comandos y tiene acceso al repositorio y al
entorno de desarrollo local.

\subsection{Usos realizados}

\subsubsection*{Apoyo al desarrollo de CloudPorter}

Claude Code se empleó como asistente de programación durante el ciclo de desarrollo iterativo
descrito en el capítulo~\ref{cap:desarrollo}. Los usos concretos incluyeron:

\begin{itemize}
    \item \textbf{Generación y revisión de código:} implementación de módulos del traductor,
    plantillas Jinja2 y la lógica del mapper determinístico, revisados y validados por el
    autor antes de cada integración.
    \item \textbf{Skills de flujo de trabajo:} se desarrollaron dos skills propias para
    Claude Code —\texttt{create-issue} y \texttt{github-workflow}— que encapsulan las
    convenciones del proyecto y asisten en la creación de issues y la revisión de workflows
    de GitHub Actions.
    \item \textbf{Reglas de contexto:} se configuró un fichero de reglas específico para
    el módulo de traducción (\texttt{.claude/rules/translator.md}) que documenta las
    convenciones de estructura del traductor y orienta al agente cuando trabaja en esa
    parte del código.
    \item \textbf{Apoyo en depuración y diseño:} consultas sobre decisiones de arquitectura,
    resolución de errores en la suite de tests y revisión de pull requests.
\end{itemize}

En todos los casos, el código generado fue revisado, adaptado y validado por el autor
mediante la suite de tests automatizados del proyecto —con un umbral mínimo de cobertura
del 80\%— y mediante la ejecución manual de los flujos de despliegue sobre AWS y Azure.

\subsubsection*{Apoyo a la redacción de la memoria}

Claude Code se utilizó como asistente de escritura académica para la redacción y revisión
de los capítulos de esta memoria. Los usos concretos incluyeron:

\begin{itemize}
    \item \textbf{Redacción de secciones:} generación de borradores de capítulos a partir
    de indicaciones del autor, respetando el enfoque narrativo del TFG y las convenciones
    de estilo definidas.
    \item \textbf{Revisión lingüística y de estilo:} corrección de coherencia, fluidez y
    adecuación al registro académico en castellano.
    \item \textbf{Skill de redacción:} se configuró una skill propia (\texttt{escritura-tfg})
    que encapsula las convenciones de voz, estructura y estilo propias de esta memoria,
    asegurando consistencia entre sesiones de trabajo.
    \item \textbf{Contexto persistente:} se mantiene un fichero \texttt{CLAUDE.md} con el
    enfoque narrativo del TFG, la estructura de capítulos y las instrucciones de redacción,
    que se carga automáticamente en cada sesión de trabajo con el agente.
\end{itemize}

Todo el contenido generado ha sido revisado críticamente por el autor, que ha verificado
la corrección técnica de las afirmaciones, la coherencia con el resto del documento y
la adecuación a los requisitos académicos. Las decisiones de contenido —qué incluir,
qué omitir, cómo enmarcar cada argumento— son en todo momento responsabilidad del autor.

\subsection{Validación}

Los mecanismos de validación empleados para garantizar la calidad del trabajo son los
siguientes:

\begin{itemize}
    \item \textbf{Código:} suite de tests automatizados con cobertura mínima del 80\%,
    comprobaciones de tipos con mypy, análisis estático con Ruff y validación funcional
    mediante despliegue real en AWS y Azure.
    \item \textbf{Memoria:} revisión manual de cada sección generada, contraste de datos
    técnicos con la documentación oficial de las herramientas y fuentes bibliográficas
    citadas, y compilación continua del documento LaTeX para verificar su integridad formal.
\end{itemize}

\subsection{Periodo de uso}

Las herramientas de IA descritas se emplearon durante todo el período de desarrollo e
implementación del proyecto, comprendido entre febrero y mayo de 2026.
